\section{Conclusioni e discussione}
\label{sec:conclusioni}

\subsection{Considerazioni sull'affidabilità del login federato}

Con il login federato l'applicazione non verifica l'utente, ma
un'attestazione emessa da qualcun altro su di lui. Finché l'attestazione è
corretta, il modello elimina i problemi visti nel Capitolo~\ref{sec:introduzione}: 
password riutilizzate, credential stuffing, database di credenziali da proteggere. 
Quando l'IdP firma un'attestazione falsa, i fattori dell'utente non contano più, 
e la firma digitale non aiuta: garantisce chi ha prodotto il token, non che il contenuto sia vero. 
Nel 2020 la verifica della firma confermava soltanto che l'errore era di Apple.

L'analisi ha dei limiti: i dettagli tecnici vengono quasi solo dal resoconto
del ricercatore, perché Apple non ha descritto il difetto~\cite{infoq2020}, e
non si sa quante applicazioni fossero vulnerabili. Gli impatti dei
Capitoli~\ref{sec:episodi} e~\ref{sec:conseguenze} sono quindi scenari
possibili, non danni documentati.

\subsection{Il caso Apple come esempio dei rischi della centralizzazione dell'autenticazione}

La centralizzazione funziona nei due sensi, come riassume la Tabella~\ref{tab:centralizzazione}.

\begin{table}[H]
    \centering
    \small
    \begin{tabular}{|p{0.22\textwidth}|p{0.33\textwidth}|p{0.33\textwidth}|}
        \hline
         & \textbf{Password gestita da ogni sito} & \textbf{Login federato} \\
        \hline
        Furto del database di un sito & Colpisce gli utenti di quel sito, e
        altri siti se la password è riutilizzata & Il sito non conserva
        password da rubare \\
        \hline
        Errore nel sistema che verifica l'identità & Riguarda un solo sito &
        Riguarda tutte le app che si affidano all'IdP \\
        \hline
        Chi corregge & Ogni sito, separatamente & L'IdP, una volta per tutti \\
        \hline
        Cosa può fare l'utente & Password unica e robusta, MFA & Poco: la
        sicurezza dipende da IdP e applicazioni \\
        \hline
    \end{tabular}
    \caption{Autenticazione locale e federata a confronto.}
    {\footnotesize Fonte: elaborazione del gruppo sui Capitoli~\ref{sec:introduzione}, \ref{sec:contesto} e~\ref{sec:mitigazioni}.\par}
    \label{tab:centralizzazione}
\end{table}

In termini di Cyber Risk Management, affidare l'autenticazione a pochi provider
riduce la probabilità di molti errori, perché ci sono meno password e meno
implementazioni da sbagliare, ma aumenta l'impatto di quelli che restano. 
Nel 2020 un solo difetto esponeva tutte le applicazioni con Sign in with Apple prive di controlli aggiuntivi, e
una sola correzione le ha protette tutte. L'esito favorevole, però, è dipeso da
chi ha trovato il difetto per primo: un ricercatore che lo ha segnalato tramite
il bug bounty~\cite{jain-apple-2020}. 
Un attaccante avrebbe potuto usarlo in silenzio, e i log sarebbero serviti a scoprirlo dopo, non a evitarlo.
 
Per chi usa questi servizi il login federato resta una scelta ragionevole
rispetto a riutilizzare la stessa password su decine di siti. 
La sua sicurezza, però, non è nelle mani dell'utente: dipende da quanto bene l'IdP controlla ciò
che firma e da quanto poco le applicazioni si fidano di un singolo claim.